% Catena biatomica: geometria, celle e integrali di sovrapposizione tra primi vicini (§2)
\begin{tikzpicture}[x=1.45cm, y=1cm]
 \tikzset{atA/.style={circle, ball color=atomA!85!white, minimum size=6.2mm, inner sep=0pt},
          atB/.style={circle, ball color=atomB!80!white, minimum size=4.2mm, inner sep=0pt}}
 \draw[black!30, thick] (-3.6,0) -- (4.6,0);
 % celle
 \foreach \c/\lab in {-2/{n-1}, 0/{n}, 2/{n+1}} {
   \draw[black!25, densely dashed] (\c-0.25,-0.75) rectangle (\c+1.25,0.62);
   \node[font=\scriptsize, black!55] at (\c+0.5,-1.02) {cella $\lab$};
 }
 % atomi
 \foreach \x in {-2,0,2} \node[atA] at (\x,0) {};
 \foreach \x in {-1,1,3} \node[atB] at (\x,0) {};
 \node[atB, opacity=0.35] at (-3,0) {};
 \node[atA, opacity=0.35] at (4,0) {};
 \node[font=\scriptsize] at (-2,-0.5) {$A_{n-1}$};
 \node[font=\scriptsize] at (-1,-0.5) {$B_{n-1}$};
 \node[font=\scriptsize] at (0,-0.5) {$A_n$};
 \node[font=\scriptsize] at (1,-0.5) {$B_n$};
 \node[font=\scriptsize] at (2,-0.5) {$A_{n+1}$};
 \node[font=\scriptsize] at (3,-0.5) {$B_{n+1}$};
 % integrali di sovrapposizione di A_n e di B_n
 \draw[-{Stealth[length=2mm]}, vA, thick] (0,0.33) to[bend right=50] node[midway, above, font=\scriptsize] {$-t$} (-1,0.26);
 \draw[-{Stealth[length=2mm]}, vA, thick] (0,0.33) to[bend left=50] node[midway, above, font=\scriptsize] {$-t$} (1,0.26);
 \draw[-{Stealth[length=2mm]}, vB, thick] (1,0.26) to[bend left=50] node[midway, above, font=\scriptsize] {$-t$} (2,0.33);
 % passo e distanza A-B
 \draw[|<->|, black!70] (0,1.15) -- (2,1.15) node[midway, above, font=\scriptsize] {$a$ (passo)};
 \draw[|<->|, black!55] (-2,1.15) -- (-1,1.15) node[midway, above, font=\scriptsize] {$a/2$};
 \node[font=\scriptsize, anchor=west] at (4.75,0) {$x$};
 \node[font=\scriptsize, black!60, anchor=north east] at (0.25,-1.35) {$x_{A_n}=na$,};
 \node[font=\scriptsize, black!60, anchor=north west] at (0.75,-1.35) {$x_{B_n}=(n+\tfrac12)a$};
\end{tikzpicture}
